% !TEX root = ../main.tex

\section{Discussion}
\label{sec:discussion}

At a fixed budget, AUF with a small limit measured mainly how far each start was from the small subset sizes, not how well the search worked among them. In every group of runs, reach accounted for 98 to 100 percent of the AUF difference between the two starts, mostly at limits that the OBL start never reached (Section~\ref{sec:res-reach}). At the sizes both starts reached, the mean difference in test F1 was small under the decision tree and LightGBM; under the random forest it was small only once the two smallest test sets were set aside, and single datasets differed by up to 0.036. A large AUF lead at a limit of ten features suggests that the leading method finds better small subsets; in the tested configurations, it meant mainly that the other method returned no small subset, or returned one in fewer runs. The same check applies to any comparison in which one method starts from, or stays near, large subsets, including comparisons of other feature-selection methods against an opposition-based NSGA-II: before a lead at a small limit is attributed to better search, it should be checked whether both methods reached that limit.

The analysis of Section~\ref{sec:metrics} and the runs give a consistent account of why this happens. Beyond a few dozen features, an OBL initial population almost never contains a subset with ten or fewer features (Section~\ref{sec:access}), so the search has to move down from about half of the features. It does so slowly: in the pilot, an OBL run first held such a subset after a median of 0.54 to 7.10 times the benchmark budget (Section~\ref{sec:res-pilot}). Once small subsets are present, the large ones disappear from the population under the tested variation and survival rules, whatever the start (Table~\ref{tab:pilot-fronts}). Small subsets then rarely change under balanced mutation, which leaves a share of about $e^{-2k/d}$ of the children unchanged, so repeated subsets dominate the proposals, and the cache keeps them out of the count of fits. The observations agree with this account at each step. We did not vary crossover or selection, however, so we do not claim that it is the only possible explanation.

The sparse start is not a better feature-selection method; it searches in a different region. Its best test F1 at any size was lower by 0.007 to 0.029 on average, and its fronts ended by 36 features (Section~\ref{sec:res-cost}). Neither tested start covered both ends of the front, and the uniform and log-uniform starts of the pilot lost their large subsets in the same way. Because the large subsets are lost during the search rather than missing at the start, keeping both ends probably requires a change in which subsets survive, for example letting the archive of the best subset of each size, which our runs keep only for reporting, take part in selection, niching by subset size, or survival by size group~\cite{Li2026CORAL}. We did not test such rules.

The improvement phase of our NSGAII-MIIP implementation did not compensate for the dense start: it raised reach by only 0.04 to 0.08, and the sparse start remained ahead of it by 0.46 to 0.64 (Section~\ref{sec:res-miip}). Within the tested budgets, the start decided whether the front reached small sizes.

For small subsets, a budget of fits is also a poor bound on running time. Without the stopping rule, single sparse runs lasted up to 563{,}390 generations in the designed experiment and 465{,}914 in the benchmark with the decision tree. The forced flip made sparse runs on the four pilot datasets about 3 to 17 times faster with the same AUF (Section~\ref{sec:res-pilot}), so it controls cost rather than improving the search. A budget of fits is a fair basis for comparing how much evaluation two methods used, but it does not measure computing time; reporting both follows the distinction between evaluation budgets and computing time in optimization benchmarking~\cite{BartzBeielstein2020}.

These results lead to a short list of quantities to report alongside AUF in a comparison at a feature limit:
\begin{itemize}
  \item whether each run reaches the limit, and after how much of the budget;
  \item the score at sizes that both methods reach, per dataset and as a mean with its interval;
  \item the best score at any size, which shows what a method gives up among large subsets;
  \item the numbers of proposed subsets and new model fits, and the elapsed time;
  \item which runs did not finish, and how they were treated in the analysis.
\end{itemize}
The budget should also be chosen with the number of features in mind. Otherwise, the comparison should be described as one of reach at that budget.

The main reading persisted under the three tested treatments of missing runs, which are scenarios rather than bounds, and under all three classifiers (Appendix~\ref{sec:appendix}).

The evidence has several limits. Most datasets come from a benchmark whose runs we analyzed after their results were known, and the pilot used datasets whose earlier results we had seen, so both guide the interpretation rather than confirm it. Thirty-eight of the 46 datasets are time series from the UCR and UEA archives, in which neighboring features are strongly correlated. Every dataset has a single split, so the five seeds measure the variability of the search, not of the data. We tested one family of operators, our implementation of those of NSGAII-MIIP, with a population of 100, and we did not compare its outputs with those of the authors' code. The benchmark budgets were calibrated on FASTENER and grow with the number of features, so the benchmark alone cannot separate the two effects. Only the pilot varied the budget, and its three budgets are points on the same search trajectories, not independent runs. Quality at matched sizes is known only for the few datasets on which both starts reach small sizes, and the rule that sets aside the smallest test sets was chosen after a first analysis. The pilot used a single classifier, the decision tree. Finally, the random split of EOData makes its test scores optimistic.
